\documentclass[aspectratio=169,19pt,english]{beamer}

% =========================================================
\usetheme{Berkeley}
\usecolortheme{default}

\definecolor{PythonBlue}{HTML}{306998}
\definecolor{PythonYellow}{HTML}{FFD43B}
\definecolor{DeepNavy}{HTML}{14213D}
\definecolor{SoftWhite}{HTML}{F7F9FC}
\definecolor{CodeGreen}{HTML}{2A9D8F}

\setbeamercolor{structure}{fg=PythonBlue}
\setbeamercolor{palette primary}{bg=DeepNavy,fg=white}
\setbeamercolor{palette secondary}{bg=PythonBlue,fg=white}
\setbeamercolor{palette tertiary}{bg=DeepNavy,fg=white}
\setbeamercolor{palette quaternary}{bg=PythonBlue,fg=white}
\setbeamercolor{title}{fg=white,bg=DeepNavy}
\setbeamercolor{frametitle}{fg=white,bg=DeepNavy}
\setbeamercolor{block title}{fg=white,bg=PythonBlue}
\setbeamercolor{block body}{fg=DeepNavy,bg=SoftWhite}
\setbeamercolor{alerted text}{fg=PythonBlue}
\setbeamercolor{normal text}{fg=DeepNavy,bg=white}
\setbeamercolor{item}{fg=PythonBlue}
\setbeamercolor{section in toc}{fg=PythonBlue}

\setbeamertemplate{navigation symbols}{}
\setbeamertemplate{items}[triangle]
\setbeamertemplate{blocks}[rounded][shadow=false]


\setbeamercovered{invisible}

% =========================================================
% PACKAGES
% =========================================================
\usepackage{enumitem} 
\usepackage[T1]{fontenc}
\usepackage[utf8]{inputenc}
\usepackage{babel}
\usepackage{lmodern}
\usepackage{microtype}
\usepackage{enumitem}
\usepackage{graphicx}
\usepackage{booktabs}
\usepackage{mathtools}
\usepackage{amssymb}
\usepackage{xcolor}
\usepackage{listings}
\usepackage{tikz}
\usepackage{hyperref}
\usepackage{listings}
% =========================================================
% PYTHON CODE STYLE
% =========================================================
\lstdefinestyle{pythonstyle}{
	language=Python,
	basicstyle=\ttfamily\small,
	keywordstyle=\color{PythonBlue}\bfseries,
	stringstyle=\color{CodeGreen},
	commentstyle=\color{gray}\itshape,
	numbers=left,
	numberstyle=\tiny\color{gray},
	numbersep=8pt,
	showstringspaces=false,
	breaklines=true,
	frame=single,
	rulecolor=\color{PythonBlue!25},
	backgroundcolor=\color{SoftWhite},
	tabsize=4,
	keepspaces=true,
	columns=fullflexible
}
\lstset{style=pythonstyle}

% =========================================================
% metadata
% =========================================================
\title[Some Basics.]{Introduction to Programming}
\subtitle{Python Fundamentals}
\author[ ]{Luis Ángel Bustos Ochoa}
\institute{Applied Physics}
\date{\today}

% =========================================================
% DOCUMENT
% =========================================================
\begin{document}
	
	% ---------------------------------------------------------
	% FRAME 1 — TITLE PAGE  ---------------------------------------------------------
	{
		\setbeamercolor{background canvas}{bg=DeepNavy}
		\begin{frame}[plain]
			\centering
			\vspace{0.7cm}
			
			{\color{PythonYellow}\rule{0.22\textwidth}{2pt}}
			
			\vspace{0.65cm}
			
			{\color{white}\Huge\bfseries Introduction To Programming\par}
			
			\vspace{0.35cm}
			
			{\color{PythonYellow}\Large Python fundamentals\par}
			
			\vspace{0.8cm}
			
			\begin{tikzpicture}
				\node[
				rounded corners=8pt,
				fill=white,
				text=PythonBlue,
				inner sep=14pt
				] {
					\begin{minipage}{0.58\textwidth}
						\centering
						\ttfamily\Large
						Universidad...
					\end{minipage}
				};
			\end{tikzpicture}
			
			\vfill
			
			{\color{white}\large Luis Ángel Bustos Ochoa\par}
			\vspace{0.15cm}
			{\color{white!75}\normalsize Some Basics \textbar{} Course: \par}
			
			\vspace{0.55cm}
			{\color{PythonYellow}\rule{0.22\textwidth}{2pt}}
			
			\vspace{0.35cm}
			{\color{white!70}\small\insertdate\par}
		\end{frame}
	}	
	%
	\section{Installation}
	{\Large
		\begin{frame}{Operations Systems}
			
			\begin{block}{WINDOWS (not preinstalled)}
				Downloading the official installer (.exe) from python.org. (It is mandatory to check the 'Add Python to PATH' box).
				
			
				
			\end{block}
			\begin{exampleblock}{MACOS(preinstalled)}
				Try to download or have a version \geq 3. 
				Download the package or archive that ends in 'pkg' in python.org.
				
				
				
			\end{exampleblock}
			
			%		{\small 
				%			\begin{tabular}{|c|c|c|}
					%				\hline
					%				WINDOWS & not preinstalled & Downloading the official installer (.exe) from python.org. (It is mandatory to check the 'Add Python to PATH' box).  \\
					%				\hline
					%				MACOS & it is pre-installated & But..  Try to download or have a version>3   \\
					%				\hline
					%				LINUX & it is pre-installated & Try 'sudo apt install python3 python3-pip python3-venv' \\
					%				\hline
					%			\end{tabular} 
				%		}
			
		\end{frame}
		
		\begin{frame}
				\begin{block}{LINUX(preinstalled)}
				Try 'sudo apt install python3 python3-pip python3-venv if you don't have a version \geq 3
				
				
			\end{block}
		\end{frame}
		
		\begin{frame}{Integrated Development Environment(IDE)}
			\framesubtitle{A complete software that make your "program experience" a lot nicer.}
			\begin{itemize}[label=$\circ$]
				\item Write your Python programs.
				\item A way to see lines of code, before running your program.
				\item A way to try and step through your code line by line.
				\item A console to interact with the user of your programs. 
			\end{itemize}
			
			
		\end{frame}
		\begin{frame}{Some Famous IDEs}
			
			\begin{itemize}[label=$\diamond$]
				\item VsCode
				\item Pycharm
				\item Spyder 
			\end{itemize}
			
		\end{frame}
		\section{Console}
		\begin{frame}{Console}
			python --version: your current python version. \\
			pip --version: the version of your package manager.
			\begin{alertblock}{TRY..}
				
			\begin{enumerate}[label=\arabic*.-]
				\item 6<7
				\item print (0)
				\item 7*0+4
				\item print ("hello world!")
			\end{enumerate}
				
			\end{alertblock}
			
		\end{frame}
		
\end{document}